% Corte mecánico íntegro: parte_ii.tex:323--419.
% Residencia material del ensamblaje sucesor de 102 capítulos.
% Residencia causal: capítulo 46.

\chapter{Tornillo espectral de Pell y suspensión logarítmica}

\section{Ley de tornillo}

El corpus contiene la identidad exacta
\begin{equation}
 V^{-1}\mathscr S V=(2+\sqrt3)\ii\,\mathscr S.
 \label{espiral:eq:tornillo-pell}
\end{equation}
Sea \(\varepsilon=2+\sqrt3\). La órbita transversal
\[
 z_n=z_0(\varepsilon\ii)^n
\]
satisface
\[
 r_n=r_0\varepsilon^n,
 \qquad
 \theta_n=\theta_0+\frac{\pi n}{2}.
\]
Eliminando \(n\),
\begin{equation}
 r(\theta)=r_0
 \exp\!\left[
 \frac{2\log(2+\sqrt3)}{\pi}
 (\theta-\theta_0)
 \right].
 \label{espiral:eq:espiral-log-pell}
\end{equation}
La espiral logarítmica es una salida del operador: no se eligió su ecuación
para ajustar después \(\varepsilon\).

\begin{proof}[Derivación de \cref{espiral:eq:espiral-log-pell}]
De \(\theta_n-\theta_0=n\pi/2\) se sigue
\(n=2(\theta_n-\theta_0)/\pi\). Sustituyendo en
\(r_n=r_0\varepsilon^n\),
\[
 r_n=r_0\exp\left(
 \frac{2\log\varepsilon}{\pi}(\theta_n-\theta_0)
 \right).
\]
La interpolación continua es la ecuación indicada.
\end{proof}

\section{Involución conjugada de contracción}

La extensión bilateral \(n\in\Z\) produce
\[
 n\to+\infty:\ r_n\to\infty,
 \qquad
 n\to-\infty:\ r_n\to0.
\]
La inversión \(n\mapsto-n\) intercambia expansión y contracción. Como
\(\varepsilon^{-1}=2-\sqrt3\), ambas ramas pertenecen a la misma unidad de
Pell. No son dos espirales independientes, sino dos orientaciones del mismo
tornillo espectral.

\begin{figure}[htbp]
\centering
\begin{tikzpicture}[scale=.82]
 \draw[->,HMTGray] (-5.1,0)--(5.1,0) node[right] {$x$};
 \draw[->,HMTGray] (0,-3.1)--(0,3.1) node[above] {$y$};
 \draw[HMTBlue,very thick,domain=-7:5.2,samples=240,smooth,variable=\t]
   plot ({.27*exp(.22*\t)*cos(deg(\t))},{.27*exp(.22*\t)*sin(deg(\t))});
 \draw[HMTRed,thick,domain=-7:5.2,samples=240,smooth,variable=\t]
   plot ({.27*exp(.22*\t)*cos(deg(-\t))},{.27*exp(.22*\t)*sin(deg(-\t))});
 \foreach \k in {-4,-3,...,4}{
   \pgfmathsetmacro{\ang}{1.57079632679*\k}
   \fill[HMTNavy]
     ({.27*exp(.22*\ang)*cos(deg(\ang))},
      {.27*exp(.22*\ang)*sin(deg(\ang))}) circle (.8pt);
 }
 \node[HMTBlue,font=\small\sffamily] at (3.5,2.4) {expansión};
 \node[HMTRed,font=\small\sffamily] at (3.5,-2.4) {orientación conjugada};
\end{tikzpicture}
\caption{Proyección transversal del tornillo de Pell. Las marcas representan cuartos de giro.}
\label{espiral:fig:tornillo-pell}
\end{figure}

\section{Covariancia de escala y conservación energética}

La dilatación \(\varepsilon\) no debe interpretarse como crecimiento
gratuito de la amplitud energética de una onda en vacío. Su realización
coherente actúa sobre la escala espacial o espectral:
\[
 (r,\varphi,\lambda)
 \longmapsto
 (\varepsilon r,\varphi+\pi/2,\varepsilon\lambda).
\]
Esta distinción será esencial en la realización electromagnética de la
\cref{espiral:chap:realizacion-em}.

\begin{resultbox}[title={Compatibilidad exacta entre el tornillo espectral de Pell y la suspensión fase--memoria}]
La elevación fase--memoria mínima es una hélice circular. El tornillo
espectral combina un cuarto de giro y una dilatación de Pell; su proyección
es una espiral logarítmica. Ambos proceden del registro HMT, pero no son la
misma periodicidad ni se obtienen el uno del otro por olvido de tipos.
\end{resultbox}
